\documentclass[12pt,twoside]{article}
%\usepackage[table]{xcolor} % important to avoid options clash.
%\input{02465shared_preamble}
%\usepackage{cleveref}
\usepackage{url}
\usepackage{graphics}
\usepackage{multicol}
\usepackage{rotate}
\usepackage{rotating}
\usepackage{booktabs}
\usepackage{hyperref}
\usepackage{pifont}
\usepackage{latexsym}
\usepackage[english]{babel}
\usepackage{epstopdf}
\usepackage{etoolbox}
\usepackage{amsmath}
\usepackage{amssymb}
\usepackage{multirow,epstopdf}
\usepackage{fancyhdr}
\usepackage{booktabs}
\usepackage{xcolor}
\newcommand\redt[1]{ {\textcolor[rgb]{0.60, 0.00, 0.00}{\textbf{ #1} } } }


\newcommand{\m}[1]{\boldsymbol{ #1}}
\newcommand{\yoursolution}{ \redt{(your solution here) } } 



\title{ Report 1 hand-in }
\date{ \today }
\author{Alen Hodzic (\texttt{s253790}) } 

\begin{document}
\maketitle

\begin{table}[ht!]
\caption{Attribution table}
\begin{tabular}{ll}
\toprule
Task & Alen \\
\midrule
1: A basic blaster-business                            & 100\% \\
2: Warmup                                              & 100\% \\
3: Manually computing $J_{N-1}$                        & 100\% \\
4: Compute optimal policy and value function           & 100\% \\
5: Kiosk2                                              & 100\% \\
6: Policy explanation continued                        & 100\% \\
7: Go east                                             & 100\% \\
8: Predict consequence of actions                      & 100\% \\
9: Possible future states                              & 100\% \\
10: Shortest path                                      & 100\% \\
11: Predict consequence of actions with one ghost      & 100\% \\
12: Possible future states with one ghost              & 100\% \\
13: Optimal one-ghost planning                         & 100\% \\
14: Predict consequence of actions with several ghosts & 100\% \\
15: Future states                                      & 100\% \\
16: Optimal planning                                   & 100\% \\
\bottomrule
\end{tabular}
\end{table}

%\paragraph{Statement about collaboration:}
%Please edit this section to reflect how you have used external resources. The following statement will in most cases suffice: 
%\emph{The code in the irls/project1 directory is entirely}

%\paragraph{Main report:}
Headings have been inserted in the document for readability. You only have to edit the part which says \yoursolution. 

\section{The kiosk (\texttt{kiosk.py})}
\subsubsection*{{\color{red}Problem 1:  A basic blaster-business}}

\yoursolution 	
\redt{Include the following elements: \begin{align}
	N & = 14 \\
	\mbox{for $k=0,\dots,N$: }\quad	\mathcal{S}_k & =  \{0,1,2,\dots,20\} \\
	\mbox{for $k=0,\dots,N-1$: }\quad \mathcal{A}_k(x_k) & = \{0,1,2,\dots,15\} \\
    \mbox{for $k=0,\dots,N-1$: }\quad f_k(x_k, u_k, w_k) & = \min\{20,\; x_k + u_k - w_k\} \\
	\mbox{for $k=0,\dots,N-1$: }\quad g_k(x_k, u_k, w_k) & = c_{o}\,u_k - s_{p}\,\min\{x_k+u_k,\; w_k\} \\
	g_N(x_N) & = 0 \\
    \mbox{for $k=0,\dots,N-1$: }\quad P_W(w_k | x_k, u_k) & = \begin{cases}
0.30 & w_k = 0 \\
0.60 & w_k = 3 \\
0.10 & w_k = 6
\end{cases}
\end{align} }

\subsubsection*{{\color{red}Problem 3:  Manually computing $J_{N-1}$}}
		
	\yoursolution The last step of the DP algorithm gives:
	$$
	J_{N-1}(20)  = \min_{u_{N-1} } \mathbb{E}_{w_{N-1}}\left[ g_{N-1}(20, u_{N-1}, w_{N-1} ) +  g_N(x_N) \right]
	$$
	Since the terminal cost is $g_N(x_N)=0$, the expression simplifies to
    $$
    J_{N-1}(20)  = \min_{u_{N-1}} \mathbb{E}_{w_{N-1}}\left[g_{N-1}(20,u_{N-1},w)\right].
    $$
    
    The expectation is taken with respect to the demand distribution
    $$
    P(w_{N-1}=0)=0.3, \quad
    P(w_{N-1}=3)=0.6, \quad
    P(w_{N-1}=6)=0.1.
    $$
    
    Therefore
    \begin{align}
    J_{N-1}(20) = \min_{u_{N-1}\in\{0,\dots,15\}}
    \Big[
    0.3\, g_{N-1}(20,u_{N-1},0) +
    0.6\, g_{N-1}(20,u_{N-1},3) +
    0.1\, g_{N-1}(20,u_{N-1},6)
    \Big].
    \end{align}
    Since the inventory is already at the maximum capacity $20$ and the maximum demand is $6$, ordering additional blasters will only increase the cost. Therefore the optimal action is
    $$
    u_{N-1}^* = 0.
    $$
    Thus
    $$
    J_{N-1}(20) = -0\cdot0.3\cdot s_{p} - 3\cdot0.6\cdot s_{p} - 6\cdot0.1\cdot s_{p} =-2.4s_{p}=-5.04.
    $$
	
\subsubsection*{{\color{red}Problem 6:  Policy explanation continued}}
	
	$$\mu_{N-1}(0) = \argmin_u \EE_{w} \left[g_k(0, u, w) + g_N( f(0,u,w) ) \right] = \cdots $$
\yoursolution
Since the terminal cost is $g_N(x_N)=0$, this simplifies to
$$
\mu_{N-1}(0) =
\operatorname*{arg\,min}_u \mathbb{E}_{w}\left[g_{N-1}(0,u,w)\right] = \operatorname*{arg\,min}_u \mathbb{E}_{w}\left[c_{o}u - s_{p}\min(u,w)\right]  .
$$

The function $\min(u,w)$ is piecewise linear in $u$ with breakpoints at
$w\in\{0,3,6\}$. Therefore the expectation $\mathbb{E}_{w}[\min(u,w)]$
is also piecewise linear in $u$, and the minimum must occur at one of
these breakpoints. Hence it suffices to evaluate $u\in\{0,3,6\}$.

\[
\begin{aligned}
u=0 &: \quad \mathbb{E}_{w}[g(0,0,w)] = 0 \\
u=3 &: \quad \mathbb{E}_{w}[g(0,3,w)] = 3c_{o}-2.1s_{p} = 0.09\\
u=6 &: \quad \mathbb{E}_{w}[g(0,6,w)] = 6c_{o}-2.4s_{p} =  3.96.
\end{aligned}
\]
Therefore

$$
\mu_{N-1}(0)=0.
$$
\section{Avoid the droid (\texttt{pacman.py)}} 
\end{document}